\section{Instrukcja obsługi}

\subsection{Struktura projektu}

Projekt składa się z kilku kluczowych katalogów i plików:

\begin{itemize}
    \item \textbf{Source/} - Katalog zawierający kody źródłowe programu głównego.
    \item \textbf{Examples/} - Przykładowe pliki wejściowe z danymi grafów.
    \item \textbf{Doc/documentation.pdf} - Plik dokumentacji projektu, w tym opis algorytmów, dowody i przykłady użycia.
    \item \textbf{Exe/multigraph\_analysis.exe} - Program będący rozwiązaniem zadania.
    \item \textbf{Makefile} - Plik automatyzujący kompilację projektu.
\end{itemize}

\subsection{Kompilacja projektu}

Do kompilacji programu i generatora grafów należy użyć polecenia:

\begin{verbatim}
make all
\end{verbatim}

Po zakończeniu kompilacji zostanie utworzony następujący plik wykonywalny:
\begin{itemize}
    \item \texttt{multigraph\_analysis.exe} - Program główny do analizy grafów.
\end{itemize}

\subsection{Uruchamianie programu głównego}

Program główny \texttt{multigraph\_analysis.exe} wymaga pliku wejściowego z opisem grafów. Plik wejściowy powinien mieć format:

\begin{verbatim}
<number_of_graphs>
<number_of_vertices_for_graph_1>
<adjacency_matrix_for_graph_1>

<number_of_vertices_for_graph_2>
<adjacency_matrix_for_graph_2>
...
\end{verbatim}

Przykład pliku wejściowego:
\begin{verbatim}
2
3
0 1 1
1 0 1
1 1 0

4
0 2 0 1
2 0 1 0
0 1 0 1
1 0 1 0
\end{verbatim}

Uruchomienie programu:
\begin{verbatim}
./multigraph_analysis.exe <input_file>
\end{verbatim}

Przykład:
\begin{verbatim}
./multigraph_analysis.exe data/input.txt
\end{verbatim}

Program przyjmuje jedynie jeden argument, jednak obsługuje wiele grafów w jednym pliku wejściowym. W celu obliczenia metryki pomiędzy grafami, należy umieścić je w jednym pliku, jeden pod drugim. Wyniki analizy grafów zostaną wyświetlone w terminalu i mogą być przekierowane do pliku wyjściowego.

\subsection{Opis funkcji programu głównego}

\begin{itemize}
    \item \texttt{handle\_arguments(int argc, char *argv[], Config *config)} - Przetwarza argumenty wiersza poleceń. Oczekuje podania pliku wejściowego.
    \item \texttt{open\_file\_with\_retry(const char *filename)} - Otwiera plik wejściowy z danymi grafów, obsługując ponowne próby w przypadku zakłóceń systemowych.
    \item \texttt{print\_cycles(GArray *output\_cycles)} - Wyświetla cykle znalezione w grafie.
    \item \texttt{analyze\_multigraph(GraphInterface *multigraph)} - Wykonuje analizę pojedynczego grafu, w tym:
    \begin{itemize}
        \item Liczenie rozmiaru grafu.
        \item Znajdowanie wszystkich cykli.
        \item Znajdowanie cykli Hamiltona.
        \item Obliczanie minimalnego rozszerzenia dla cyklu Hamiltona.
    \end{itemize}
\end{itemize}
